\section{Reproduction by the original real Gaussian kernel}

Use the frame domain and the Lebesgue measure in ICL14. Write
$p=X.p_0$, $q=r-p$, $J=\operatorname{diag}(1_p,-1_q)$, and
$S_\delta$ for the open slab on which the Gram matrix at each CM
embedding differs entrywise from $J$ by less than $\delta$. Fix the
original finite family $P_b$ of harmonic row-highest polynomials of
the indicated weights, with at least one nonzero member. Write $K$ for
its actual doubly differentiated Gaussian kernel and $\lambda$ for
its Borel weight. We prove (ICL.7) on the entire domain. The
archimedean calculation needs even $n\ge 2r$, so includes the range
of ICL14.

% BEGIN K74A1
\paragraph{K74A1. The integrable determinant on a split Gram slab.}
At one place use ordinary complex matrix coordinates
$Y=\left(\begin{smallmatrix}A&B\\C&D\end{smallmatrix}\right)$
with $A$ of size $p\times p$. Entrywise conjugation interchanges the
source convention $Y^tJ\bar Y$ with $Y^*JY$ without altering volumes
or determinant norms. We claim that for sufficiently small $\delta$
the function $|\det A|^{-n}$ is integrable over the full slab
$Y^*JY\simeq J$ (here $\simeq$ means an entrywise error $<\delta$).

Put $U=\binom AC$, $G=U^*JU$. By shrinking $\delta$,
$G$ is uniformly near $I$ in operator norm, say
$\frac34 I<G<\frac54 I$. Conditional on $U$, each remaining column
$v$ satisfies bounds on $U^*Jv$ and $-v^*Jv$ by constants depending
only on $r$. Consequently $v$ lies in an ellipsoid
$v^*Hv\le R$ with $H=2JU U^*J-J$ and $R$ uniform. Indeed $H$ is
positive on the $J$-negative orthogonal complement of the columns
of $U$, where it equals $-J$. On the positive column space its
form is $2G^2-G$, which is positive. The spaces are orthogonal also
for $H$. The matrix determinant lemma gives
$|\det H|=|\det(I-2G)|$, uniformly bounded away from zero.
The real volume of this complex ellipsoid is its ball volume
divided by $\det H$. Dropping all cross-conditions among the last
$q$ columns bounds their joint volume by a uniform constant.

For the remaining $A,C$, positivity of $G=A^*A-C^*C$ makes $A$
invertible. Substitute $C=ZA$. The Jacobian is $|\det A|^{2q}$.
Set $P=I-Z^*Z$, which is positive since $G=A^*PA>0$. Thus $Z$
ranges in a bounded matrix ball. For fixed $Z$ set
$A'=P^{1/2}A$, so $G=A'^*A'$, and $A'$ ranges in a bounded set
with determinant bounded away from zero. The factors involving $P$
in the volume times $|\det A|^{-n}$ now multiply to
\[
                 (\det P)^{(n-2q)/2-p}.
\]
This is at most $1$ since $0<P\le I$ and $n\ge2(p+q)$. Integration
over the two bounded sets proves the claim, including the endpoint
$n=2r$. Empty matrix blocks give the same formulas with empty
determinant $1$. Left multiplication by an invertible fixed complex
matrix and permutation of coordinates change Lebesgue measure by a
constant positive factor. The claim therefore applies in any of the
original frame coordinates.
% END K74A1

% BEGIN K74A2
\paragraph{K74A2. Uniqueness of the lowest weight line.}
Fix an orthonormal $x_*$. For the product of unitary diagonal tori
let $a_t=x_*t x_*^{-1}$, $\chi(t)=\prod_{\tau,j}t_{\tau j}^{
\lambda_{\tau j}}$, with normalized Haar measure $dt$.
Two holomorphic functions of the given right weight and this left
torus character satisfy $G_1(y)G_2(x_*)=G_2(y)G_1(x_*)$. To prove
it, subtract the two sides and call the difference $m$. On the big
cell write $y=x_*Lb$, with $L$ lower triangular, $b$ upper triangular,
and use right covariance. Let
$D(s)=\operatorname{diag}(1,s,\ldots,s^{r-1})$ at each embedding.
Then $D(s)LD(s)^{-1}$ has entries $s^{i-j}L_{ij}$ for $i\ge j$,
and extends to the diagonal of $L$ at zero. The resulting frame
stays in the domain for $|s|\le1$. In the signature test on each
positive partial column space, the rows beyond the signature break
are multiplied by powers of $|s|$ no larger than those on the
positive rows. Hence their negative contribution cannot destroy
positive definiteness. The matrix stays invertible for $s\ne0$. Together
with positivity of the first $p$ columns this gives the original
domain condition. Right diagonal rescaling preserves it, and at zero
the lower-triangular factor becomes an invertible diagonal. On $|s|=1$ the left and right diagonal
characters cancel, so $m(x_*D(s)LD(s)^{-1})=m(x_*L)$.
The function is holomorphic on the disk and continuous on its
boundary. It is constant, and its center value is zero by right
covariance and $m(x_*)=0$. This proves vanishing on the big cell.
Frames with nonzero leading minors are dense in the open domain
(perturb by a small scalar multiple of $x_*$), so $m=0$ throughout.
For a function with nonzero value at the orthonormal frame, the
same calculation gives its normalized value $1$ on the lower
unipotent cell.
% END K74A2

% BEGIN K74A3
\paragraph{K74A3. The actual Gaussian, its weight and positive diagonal.}
At frames $x,y$ the kernel is obtained by applying the constant
coefficient operators $P_b$ in $z$ and $\overline{P_b}$ in $w$
to the doubled Gaussian at $z=w=0$, and summing over $b$. In the
doubled-frame expression, split the quadratic exponent into its
$z,z$, $w,w$ and $z,w$ blocks. Within each of the first two blocks
the only quadratic factors are sums of the form
$\sum_a Z^+_{ia}Z^-_{ja}$ (and their $w$ counterparts), with
coefficients independent of the column $a$. This follows by writing
the frame in its positive and negative row blocks: the same-sign
diagonal blocks of the pure quadratic correction vanish, and the
off-diagonal blocks pair the two row types. In the derivative pairing,
multiplication by this quadratic is adjoint to
$\sum_a\partial^+_{ia}\partial^-_{ja}$. Harmonicity annihilates
every positive power of each pure correction, so these exponentials
can be removed in computing the derivative. The mixed exponential
and its determinant prefactor are holomorphic in the first frame.
Indeed their entries are obtained by inverting the transverse
mixed frame matrices, whose determinants do not vanish on the
domain. Right triangular multiplication substitutes its positive
and negative blocks in the two sets of letters. The highest-weight
equations for $P_b$ and the determinant factor give precisely
$K(xb,y)=\operatorname{bchar}_\lambda(b)K(x,y)$.
Interchanging $x,y$ conjugates the mixed Gaussian and exchanges the
two derivative operators. An isometry of the Hermitian space on
the left of both frames preserves their mixed Gram expressions,
and its determinant has modulus one. Hence
$K(y,x)=\overline{K(x,y)}$ and $K(ux,uy)=K(x,y)$.

For an orthonormal frame $x_*$ the mixed exponential is
$\exp(\sum_l t_l z_l w_l)$ in the dual letter coordinates, with
$t_l=-\pi\,\mathrm{Im}(\tau_l\kappa_0)>0$. The prefactor is
$D_* =\prod_\tau |\det x_{*,\tau}|^n>0$. Expanding the
exponential and differentiating gives
\[
      K(x_*,x_*)=
        D_*\sum_{b,\alpha}\alpha!\,
           |(P_b)_\alpha|^2\prod_l t_l^{\alpha_l}>0 .
\]
This proves positivity without assuming that a chosen theta section
has nonzero value at any frame.

For later use we describe why $K(x,\cdot)$ is both bounded and
integrable on all sufficiently thin $S_\delta$, with $x$ fixed.
First take $x=x_*$ and an orthonormal $y$. In coordinates
$Y=x_*^{-1}y$ the holomorphic function $K(y,x_*)$ has right Borel
weight $\lambda$ and the corresponding left torus weight, by
invariance and symmetry above. Normalize its value at the identity.
On the big cell its value is the product of the leading minors
$D_j(Y)$ with exponents $\lambda_j-\lambda_{j+1}$, including the
last determinant with exponent $\lambda_r$. Indeed on the lower
unipotent cell its normalized value is $1$ by the contraction of
K74A2, and the upper triangular diagonal in Gaussian elimination
is $D_j/D_{j-1}$. Equalities extend to bounds at all orthonormal
frames by continuity. For $Y^*JY=J$, $|D_p|\ge1$,
$|D_r|=1$ and $|D_j|\le |D_p|$ for all $j$. To see the last
inequality, exchange the last $q$ output coordinates with the
last $q$ input coordinates in the isometry $Y$. Conservation of
the indefinite form gives a unitary map on the resulting positive
definite coordinates. Its appropriate principal minors are
$D_j/D_p$ up to sign or conjugation, by block elimination, so their
norms are at most one. Equivalently for $j\le p$, the upper $p$
block has all singular values at least one and its $j$-minors are
bounded by $|D_p|$. For $j>p$ apply the complementary-minor formula
and the same reasoning to the lower $q$ block of $Y^{-1}$.

The weights are decreasing within either signature block, and across
the break $\lambda_i+n\le\lambda_j$ for $i\le p<j$.
Bounding all the nonnegative minor exponents and retaining the
one at the break therefore gives $|K(y,x_*)|\le
C_*\prod_\tau |D_p(Y_\tau)|^{-n}$. In the definite cases the
orthonormal set is compact and this bound also holds. On the thin
slab signed Gram--Schmidt writes $y=y' b$ with $y'$ orthonormal
and the diagonal entries of $b$ and their inverses uniformly
bounded: the Gram minors are close to those of $J$ and their
successive ratios give the squared diagonal moduli. Right Borel
covariance and $D_p(Y'b)=D_p(Y')\prod_{j\le p} b_{jj}$ extend
both the bound and this integrable majorant to the slab. K74A1
applies at each of the finitely many places and Fubini gives
integrability on the product. Every fixed $x$ is an orthonormal
frame times a right Borel factor, and left unitary transport preserves
the slab and Lebesgue measure. Both assertions follow for this $x$.
% END K74A3

% BEGIN K74A4
\paragraph{K74A4. The reproducing scalar on the original slab.}
Choose $\delta_1>0$ small enough for K74A3. For every
$0<\delta\le\delta_1$, choose an orthonormal $x_*$. Put
\[
   \rho_\delta=
   \frac{\int_{S_\delta}|K(x_*,y)|^2\,dy}{K(x_*,x_*)}.
\]
The denominator is positive by K74A3. The numerator is finite by
the simultaneous integrability and boundedness there. It is
positive because the continuous integrand is positive at $x_*$
and the slab is an open neighborhood of that point. These arguments
include a zero-dimensional frame space, with its ordinary Haar
measure.

For a bounded
holomorphic form $F$ of the given right weight set
\[
                  G(y)=\int \overline{\chi(t)}F(a_ty)\,dt.
\]
The map $G$ is holomorphic: on compact subneighborhoods the
translated arguments range over a compact subset of the domain, so
local uniform bounds justify integration of the holomorphic
functions. It has the same right weight, left torus character
$\chi$, and $G(x_*)=F(x_*)$ by right covariance. For fixed
$k(y)=K(x_*,y)$ one has $k(a_ty)=\chi(t)^{-1}k(y)$.
The $a_t$ preserve both the slab and Lebesgue measure (their complex
determinants are units of modulus one). A form bounded in its
weighted norm is bounded on the slab by the Borel diagonal bounds
of K74A3. Thus Fubini and a change of variables give
$\int_{S_\delta} kG=\int_{S_\delta} kF$.

The function $y\mapsto K(y,x_*)$ has the two transformation laws of
K74A2 by K74A3. Apply uniqueness there with the torus mean $G$:
\[
   G(y) K(x_*,x_*)=K(y,x_*)F(x_*).
\]
Multiply by $K(x_*,y)$ and integrate. Symmetry of $K$ and the
torus-mean identity give the reproducing formula at $x_*$. For
any orthonormal $x'=u x_*$ apply this formula to $F(u\,\cdot)$.
Left multiplication by $u$ preserves the kernel, slab and measure.
For an arbitrary domain frame $x=x'b$ apply the right Borel law to
both sides, with the same factor. Thus for every bounded holomorphic
form with the specified Borel weight and every domain frame,
\[
             \int_{S_\delta} K(x,y)F(y)\,dy
                          =\rho_\delta F(x).
\]
The same strictly positive scalar works for all $F,x$ at this
width. The slab, measure, and kernel in this formula are those of
(ICL.7). The separate arithmetic identity (ICL.6) involves a sum
over the rational unitary group, a class sum, and the original
vacuum multipliers. That identity and the geometric level-system
construction remain necessary to apply ICL14.
% END K74A4
